% ============================================================================
\chapter{Дискретная экспонента Эйлера на \texorpdfstring{$\mathbb{Z}_{512}$}{Z512}, однородная проективная арифметика над \texorpdfstring{$\mathbb{Q}(i,\sqrt{2})$}{Q(i, sqrt 2)} и выход в непрерывный анализ}
\label{ch:mono-rotator-calculus}
% ============================================================================

\section{Дискретная экспонента Эйлера, код Грея, соотношения Коши--Римана и тождество постоянного вронскиана}

На дискретном кольце $\mathbb{Z}_{512}$ гармонические функции континуума $\cos(2\pi\varphi/512)$ и $\sin(2\pi\varphi/512)$ принимают иррациональные значения и не могут вычисляться внутри 9-битного регистра без потери обратимости. Фундаментальный результат дискретного анализа на носителе состоит в том, что существует \textbf{единственная} кусочно-линейная целочисленная пара проекций $(\mathcal{C}_3(\varphi), \mathcal{S}_3(\varphi)) \in [-128, +128]^2$, сохраняющая все дифференциальные, когомологические, голоморфные (дискретные условия Коши--Римана из \texttt{MODEL}) и симплектические (унитарное сохранение фазовой площади) инварианты тригонометрии в точных целых числах.

Представим любую 9-битную фазу $\varphi \in \{0, \dots, 511\}$ в виде $\varphi = 256\,b_8 + 128\,b_7 + m_7$, где $b_8, b_7 \in \{0,1\}$ --- два старших бита квадранта, а $m_7 \in \{0, \dots, 127\}$ --- 7-битная мантисса внутри четверти кольца $Q_\square = 128$. Введём комплементарную мантиссу $\overline{m}_7 = 128 - m_7$. Тогда чётная (косинусная) и нечётная (синусная) проекции дискретного характера Эйлера определяются одной транспозицией бабочки по 2-битному коду Грея $(g_1, g_0) = (b_8, b_8 \oplus b_7)$:
\begin{align}
  \bigl(\mathcal{C}_3(\varphi),\ \mathcal{S}_3(\varphi)\bigr) &= (+\overline{m}_7,\ +m_7) \qquad \text{при } (b_8, b_7) = (0, 0) \text{ (квадрант I, Грей } 00\text{)}, \label{eq:mono-rotator-def} \\
  \bigl(\mathcal{C}_3(\varphi),\ \mathcal{S}_3(\varphi)\bigr) &= (-m_7,\ +\overline{m}_7) \qquad \text{при } (b_8, b_7) = (0, 1) \text{ (квадрант II, Грей } 01\text{)}, \nonumber \\
  \bigl(\mathcal{C}_3(\varphi),\ \mathcal{S}_3(\varphi)\bigr) &= (-\overline{m}_7,\ -m_7) \qquad \text{при } (b_8, b_7) = (1, 0) \text{ (квадрант III, Грей } 11\text{)}, \nonumber \\
  \bigl(\mathcal{C}_3(\varphi),\ \mathcal{S}_3(\varphi)\bigr) &= (+m_7,\ -\overline{m}_7) \qquad \text{при } (b_8, b_7) = (1, 1) \text{ (квадрант IV, Грей } 10\text{)}. \nonumber
\end{align}

\begin{theorem}[Дискретные условия Коши--Римана, постоянный разностный вронскиан и замкнутые дискретные интегралы]
\label{thm:mono-wronskian-integrals}
На всех $512$ состояниях кольца $\varphi \in \mathbb{Z}_{512}$ операторы~\eqref{eq:mono-rotator-def} удовлетворяют следующим точным тождествам:
\begin{enumerate}[label=\textbf{\arabic*.},noitemsep]
\item \textbf{$L^1$-норма грани кубооктаэдра, поворот на $i$ (четверть кольца $+Q_\square = +128$) и дискретные соотношения Коши--Римана:}
  \begin{align}
    |\mathcal{C}_3(\varphi)| + |\mathcal{S}_3(\varphi)| &\equiv Q_\square(3) \;=\; 128, \label{eq:mono-l1-norm} \\
    \mathcal{C}_3(\varphi + 128) &= -\mathcal{S}_3(\varphi), \label{eq:mono-rotator-shift-c} \\
    \mathcal{S}_3(\varphi + 128) &= +\mathcal{C}_3(\varphi), \label{eq:mono-rotator-shift-s}
  \end{align}
  причём внутри каждого квадранта прямые разности вдоль кольца и по нормали к нему удовлетворяют сопряжённым дискретным условиям Коши--Римана $\Delta\mathcal{C}_3(\varphi) = -\mathrm{sgn}\bigl(\mathcal{S}_3(\varphi + 1/2)\bigr)$ и $\Delta\mathcal{S}_3(\varphi) = +\mathrm{sgn}\bigl(\mathcal{C}_3(\varphi + 1/2)\bigr)$ ($\Delta\mathcal{C}_3(\varphi + 128) = -\Delta\mathcal{S}_3(\varphi)$, $\Delta\mathcal{S}_3(\varphi + 128) = +\Delta\mathcal{C}_3(\varphi)$).
\item \textbf{Тождественное постоянство разностного вронскиана (дискретная унитарность и отсутствие точек ветвления на всём кольце):}
  для прямых разностей $\Delta\mathcal{C}_3(\varphi) = \mathcal{C}_3(\varphi+1) - \mathcal{C}_3(\varphi)$ и $\Delta\mathcal{S}_3(\varphi) = \mathcal{S}_3(\varphi+1) - \mathcal{S}_3(\varphi)$ дискретный вронскиан тождественно равен четверти кольца:
  \begin{equation}
    \label{eq:mono-wronskian}
    \mathcal{W}_3(\varphi) \;:=\; \mathcal{C}_3(\varphi)\,\Delta\mathcal{S}_3(\varphi) - \mathcal{S}_3(\varphi)\,\Delta\mathcal{C}_3(\varphi) \;\equiv\; 128 \;=\; Q_\square(3) \qquad \forall\,\varphi \in \mathbb{Z}_{512}.
  \end{equation}
\item \textbf{Замкнутая интегрируемость дискретных кольцевых сумм типа Френеля, Якоби, Дирихле, Парсеваля и Бесселя в~$\mathbb{Z}$:}
  на стационарной параболе Верле $\varphi(t) = t^2 \bmod 512$ (имеющей первый стационарный возврат на такте $T_{\mathrm{ret}} - 1 = 1023 = 2^{10}-1$ и полный двухвитковый спинорный период $2\phaseRingCount = 1024$) дискретные кольцевые суммы типа Френеля, эллиптические суммы типа Якоби на лемнискатном модуле $k^2 = 1/(c_0/c)^2 = 1/2$, сумма типа Дирихле, сумма энергии Парсеваля и сумма типа Бесселя нулевого порядка вычисляются точно в целых числах:
  \begin{align}
    \mathcal{I}_{\mathrm{Fresnel},C}(512) &= \sum_{\varphi=0}^{511} \mathcal{C}_3(\varphi^2) \;=\; \sum_{\varphi=0}^{511} \mathcal{S}_3(\varphi^2) \;=\; 2\,048, \label{eq:mono-closed-integrals} \\
    6\sum_{\varphi=0}^{511} \mathcal{C}_3(\varphi)^2 &= 6\sum_{\varphi=0}^{511} \mathcal{S}_3(\varphi)^2 \;=\; 512\,(2\cdot 128^2 + 1) \;=\; 16\,777\,728, \label{eq:mono-parseval-sum} \\
    \bigl(\mathcal{J}_0(0),\ \mathcal{J}_0(1),\ \mathcal{J}_0(4)\bigr) &= (65\,536,\ 32\,768,\ 0). \label{eq:mono-bessel-sum}
  \end{align}
\end{enumerate}
\emph{Формальная верификация:} теоремы \texttt{thm1\_\allowbreak constant\_\allowbreak wronskian\_\allowbreak weier\allowbreak strass\_\allowbreak lift\_\allowbreak and\_\allowbreak newton\_\allowbreak leibniz} и \texttt{thm2\_\allowbreak closed\_\allowbreak form\_\allowbreak fresnel\_\allowbreak jacobi\_\allowbreak dirichlet\_\allowbreak parseval\_\allowbreak bessel\_\allowbreak integrals} в модулях \texttt{papers/\allowbreak Rotator\allowbreak Quadrature.lean} и \texttt{papers/\allowbreak Logic\allowbreak Cohomology.lean}.
\end{theorem}

\begin{proof}
В пункте~(1) внутри каждого из четырёх квадрантов с локальной мантиссой $m_7 \in [0, 128]$ проекции равны $\pm(128 - m_7)$ и $\pm m_7$, сумма модулей которых тождественно равна $(128 - m_7) + m_7 = 128$. Сдвиг фазы на $+128$ переводит квадрант $(b_8, b_7)$ в следующий по коду Грея, заменяя $(\mathcal{C}_3, \mathcal{S}_3) \mapsto (-\mathcal{S}_3, +\mathcal{C}_3)$.

В пункте~(2) внутри первого квадранта ($0 \le m < 128$) соседние значения равны $(128 - m, m)$ и $(127 - m, m + 1)$, а в угловой вершине $m = 128$, где наклон $\Delta\mathcal{C}_3$ меняет знак с $-1$ на $+1$, координата $\mathcal{C}_3(128) = 0$ обращается в нуль:
\begin{align*}
  \mathcal{W}_3(m) &= (128 - m)(m + 1) - m(127 - m) \;=\; 128 \qquad (0 \le m < 128), \\
  \mathcal{W}_3(128) &= 0 \cdot 127 - 128 \cdot (-1) \;=\; 128;
\end{align*}
по симметрии четверти оборота то же равенство выполняется во всех четырёх вершинах алмаза.

В пункте~(3) для суммы Френеля~\eqref{eq:mono-closed-integrals} разбиение $512$ вычетов $\varphi \in \{0,\dots,511\}$ по степени делимости на~$2$ показывает, что $256$ нечётных точек, $192$ точки $\varphi \equiv 2, 4, 6 \pmod 8$ и $32$ точки $\varphi \equiv 0 \pmod{16}$ попарно сокращаются при сдвиге квадратичной фазы на полупериод $+256$ ($\mathcal{E}_3(\psi + 256) = -\mathcal{E}_3(\psi)$), тогда как $32$ точки $\varphi = 8(2k + 1)$ дают $\varphi^2 \equiv 64 \pmod{512}$ на биссектрисе стационарной фазы $\psi = 64$, где $\mathcal{C}_3(64) = \mathcal{S}_3(64) = 64$, откуда сумма равна $32 \times 64 = 2\,048$. Для суммы Парсеваля~\eqref{eq:mono-parseval-sum} сумма квадратов по четырём квадрантам длины $Q = 128$ равна $2 Q^2 + 4\sum_{m=1}^{Q-1} m^2 = 4Q(2Q^2 + 1)/6$, откуда $6\sum \mathcal{C}_3(\varphi)^2 = 512(2\cdot 128^2 + 1) = 16\,777\,728$. Для суммы Бесселя~\eqref{eq:mono-bessel-sum} при $k = 0$ имеем $512\times 128 = 65\,536$, при $k = 1$ среднее значение $128 - |\mathcal{S}_3(\varphi)|$ равно $64$ ($512\times 64 = 32\,768$), а при $k = 4$ аргумент $4\mathcal{S}_3(\varphi)$ равномерно пробегает полные периоды кольца, давая $\mathcal{J}_0(4) = 0$.
\end{proof}

\section{Вывод классического интеграла Римана и формулы Ньютона--Лейбница в макропределе \texorpdfstring{$N \to \infty$}{N -> infty}}

Покажем, что классический континуальный математический анализ не противоречит дискретному носителю, а возникает из него как точная макроскопическая асимптотика при суммировании по архимедовым оболочкам кубооктаэдра.

\begin{theorem}[Рождение непрерывного интеграла Римана из слоёв кубооктаэдра $S_k(3) = 10k^2 + 2$]
\label{thm:mono-riemann-emergence}
Пусть $V_{O_2}(3, N) = 1 + \sum_{k=1}^{N} (10k^2 + 2)$ --- полное число узлов в $N$-слойном причинном шаре кубооктаэдра. Тогда:
\begin{enumerate}[label=\textbf{\arabic*.},noitemsep]
\item Точный целочисленный объём шара удовлетворяет полиномиальному тождеству третьей степени без остатка:
  \begin{align}
    3\,V_{O_2}(3, N) &= 10\,N^3 + 15\,N^2 + 11\,N + 3, \label{eq:mono-ball-polynomial} \\
    \bigl(V_{O_2}(3,1),\ V_{O_2}(3,2),\ V_{O_2}(3,3)\bigr) &= (13,\ 55,\ 147). \label{eq:mono-ball-123}
  \end{align}
\item При макроскопическом масштабировании шага решётки $h = R / N \to 0$ нормированная сумма по сферическим слоям кубооктаэдра сходится к непрерывному радиальному интегралу Римана $\int_0^R r^2\,\mathrm{d}r = R^3 / 3$ с точным главным коэффициентом $|N(3)| - (c_0/c)^2 = 12 - 2 = 10$:
  \begin{equation}
    \label{eq:mono-archimedes-limit}
    \lim_{N \to \infty} \frac{1}{10\,N^3} \sum_{k=1}^{N} S_k(3) \;=\; \lim_{N \to \infty} \frac{V_{O_2}(3, N) - 1}{10\,N^3} \;=\; \frac{1}{3} \;=\; \int_0^1 x^2\,\mathrm{d}x.
  \end{equation}
\item Для любой дискретной первообразной $F(\Phi) = \sum_{\varphi=0}^{\Phi-1} f(\varphi)$ на кольце $\mathbb{Z}_{512}$ выполняется точная формула Ньютона--Лейбница без ошибки округления:
  \begin{align}
    \sum_{\varphi=a}^{b-1} \Delta F(\varphi) &= F(b) - F(a), \label{eq:mono-newton-leibniz} \\
    \oint_{\mathbb{Z}_{512}} \Delta F(\varphi) &= 0. \label{eq:mono-newton-leibniz-loop}
  \end{align}
\end{enumerate}
\emph{Формальная верификация:} теоремы \texttt{thm3\_\allowbreak continuum\_\allowbreak integration\_\allowbreak from\_\allowbreak cuboctahedron\_\allowbreak archimedes\_\allowbreak layers} и \texttt{thm5\_\allowbreak constructive\_\allowbreak de\_\allowbreak rham\_\allowbreak cohomology\_\allowbreak nilpotent\_\allowbreak algebra\_\allowbreak and\_\allowbreak banach\_\allowbreak tarski\_\allowbreak free\_\allowbreak measure} в модулях \texttt{papers/\allowbreak Rotator\allowbreak Quadrature.lean} и \texttt{papers/\allowbreak Logic\allowbreak Cohomology.lean}.
\end{theorem}

\begin{proof}
Подставляя формулу суммы квадратов $6\sum_{k=1}^N k^2 = N(N+1)(2N+1) = 2N^3 + 3N^2 + N$ в определение $3 V_{O_2}(3, N) = 3 + 5\cdot 6\sum_{k=1}^N k^2 + 6N$, получаем $3 + 5(2N^3 + 3N^2 + N) + 6N = 10N^3 + 15N^2 + 11N + 3$, что доказывает~\eqref{eq:mono-ball-polynomial}. Деление $(V_{O_2}(3,N) - 1)$ на $10N^3$ даёт $\frac{1}{3} + \frac{1}{2N} + \frac{11}{30N^2} \to \frac{1}{3}$ при $N \to \infty$, доказывая~\eqref{eq:mono-archimedes-limit}. Равенство~\eqref{eq:mono-newton-leibniz} следует из телескопического сокращения последовательных разностей $\Delta F(\varphi) = F(\varphi+1) - F(\varphi)$.
\end{proof}

\section{Однородная проективная арифметика полей, башня производных отношения и многомерная кубатура Симпсона}

\begin{theorem}[Однородная проективная арифметика в $\mathbb{Q}(i)$, $\mathbb{Q}(\sqrt{d-1})$ и $\mathbb{Q}(i,\sqrt{2}) = \mathbb{Q}(\zeta_8)$, башня разностных производных всех порядков и многомерная кубатура Симпсона на $O_2(d)$]
\label{thm:mono-projective-fields-c-infinity-simpson}
В модуле \texttt{papers/RotatorQuadrature.lean} (\texttt{thm34}, \texttt{thm35}, \texttt{thm38}, \texttt{thm39}) установлены следующие алгебраические и аналитические законы:
\begin{enumerate}[label=\textbf{(\arabic*)},leftmargin=*]
\item \textbf{Отсутствие делителей нуля на $L^1$-алмазе и проективная обратимость в полях $\mathbb{Q}(i)$, $\mathbb{Q}(\sqrt{m})$ ($m = d - 1$) и $\mathbb{Q}(i,\sqrt{2}) = \mathbb{Q}(\zeta_8)$:}
  на $L^1$-алмазе $|C| + |S| = Q \ge 1$ знаменатель Вейерштрасса $R_T = C^2 + S^2$ строго положителен, а однородные обратные в $\mathbb{Q}(i)$ и $\mathbb{Q}(\sqrt{m})$ задаются целочисленными тройками без деления:
  \begin{align}
    2(C^2 + S^2) &= (C + S)^2 + (C - S)^2 \;\ge\; Q^2 \;\ge\; 1, \label{eq:mono-diamond-no-zero-div} \\
    \operatorname{gaussFieldInv}(X, Y, R) &= (R X,\; -R Y,\; X^2 + Y^2), \label{eq:mono-gauss-inv} \\
    \operatorname{quadFieldInv}_m(a, b, D) &= (D a,\; -D b,\; a^2 - m b^2). \label{eq:mono-quad-inv}
  \end{align}
  При $d = 3$ ($m = 2$) первообразный корень $8$-й степени из единицы $\zeta_8 = (\sqrt{2} + i\sqrt{2})/2 \leftrightarrow (0, 1, 0, 1, 2) \in \mathbb{Q}(i, \sqrt{2})$ совпадает с биссектрисой стационарной фазы Френеля $\psi = K_P(3) = 64$ и удовлетворяет $\zeta_8^2 \sim i$, $\zeta_8^4 \sim -1$, $\zeta_8^8 \sim +1$, $\zeta_8\bar{\zeta}_8 \sim +1$.
\item \textbf{Замкнутая башня производных отношения $D^n(f/g) = N_n(x)/G_n(x)$ во всех порядках $n \in \mathbb{N}$, безусловная обратимость знаменателя Вейерштрасса и вариационная производная Эйлера--Лагранжа:}
  для оператора $(Df)(x) = f(x+1) - f(x)$ и отношения $f / g$ индуктивный числитель $N_n(x)$ и знаменатель $G_n(x)$ задаются рекуррентными формулами, причём для однородного знаменателя Вейерштрасса $g(t) = R_3(t) = \mathcal{C}_3(t)^2 + \mathcal{S}_3(t)^2$ безусловно (при всех $n \in \mathbb{N}$ и всех $x \in \mathbb{Z}$) выполняется $G_n(x) \ge 1$:
  \begin{align}
    N_0(x) &= f(x), \qquad
    N_{n+1}(x) \;=\; N_n(x+1)\,g(x) - N_n(x)\,g(x+n+1), \label{eq:mono-quotient-numer-rec} \\
    G_n(x) &= \prod_{j=0}^n g(x+j) \;\ge\; 1, \label{eq:mono-quotient-denom-prod} \\
    \bigl(N_n(x+1)\,G_n(x) &- N_n(x)\,G_n(x+1)\bigr)\,g(x) \;=\; N_{n+1}(x)\,G_n(x). \label{eq:mono-quotient-deriv-step}
  \end{align}
  На составной траектории $(L_1, w, L_2)$ функционал действия $\mathcal{S}(D_w) = Q_2(Q_w D_1 + P_1 D_w) + P_1 P_w D_2$ имеет первую вариацию
  \begin{equation}
    \label{eq:mono-action-variation}
    \delta\mathcal{S} \;=\; P_1\,(D_w' - D_w)\,Q_2 \qquad \bigl(\delta\mathcal{S} = 0 \iff D_w' = D_w \text{ при } P_1, Q_2 \ge 1\bigr).
  \end{equation}
\item \textbf{Многомерное разбиение на пояса и кубатура Симпсона на $O_2(d)$ при произвольной размерности $d \ge 2$:}
  на оболочке $O_2(d) = \{e \in \{-1,0,+1\}^d \mid \sum_{i=1}^d e_i^2 = 2\}$ из $|V_{O_2}(d)| = 2d(d-1)$ вершин сечение вдоль любой оси образует три пояса, которые в сумме с центральным весом $w_0(d) = 2(d-1)(6-d)$ дают параболические веса Симпсона:
  \begin{align}
    \bigl(|Z_{-1}(d)|,\ |Z_0(d)|,\ |Z_{+1}(d)|\bigr) &= \bigl(2(d-1),\ 2(d-1)(d-2),\ 2(d-1)\bigr), \label{eq:mono-belt-partition-d} \\
    \bigl(|Z_{-1}(d)|,\ |Z_0(d)| + w_0(d),\ |Z_{+1}(d)|\bigr) &= 2(d-1) \times (1,\ 4,\ 1), \label{eq:mono-simpson-weights-d}
  \end{align}
  становясь равновеликими $(4, 4, 4)$ с $w_0(3) = 12 = |V_{O_2}(3)|$ при $d = 3$ и воспроизводя точный трёхмерный интеграл Римана $3\iiint_{[-1,1]^3} P\,\mathrm{d}V = 24c_0 + 8(c_x+c_y+c_z)$ на любом кубическом многочлене $P$.
\end{enumerate}
\end{theorem}

\begin{proof}
В пункте~(1) при $|C| + |S| = Q \ge 1$ одно из чисел $|C + S|$ или $|C - S|$ равно $Q$, откуда $(C + S)^2 + (C - S)^2 \ge Q^2 \ge 1$, а перемножение $(a, b, D)\otimes_{\sqrt{m}}(Da, -Db, a^2 - m b^2) = (D(a^2 - m b^2), 0, D(a^2 - m b^2)) \sim (1, 0, 1)$ проверяется раскрытием скобок. В пункте~(2) индукцией по $n \in \mathbb{N}$ доказывается $G_n(x+1)g(x) = g(x+n+1)G_n(x) = G_{n+1}(x)$, подстановка которого в левую часть~\eqref{eq:mono-quotient-deriv-step} выносит общий множитель $G_n(x)$. В пункте~(3) раскрытие $2(d-1)(d-2) + 2(d-1)(6-d) = 2(d-1)(d - 2 + 6 - d) = 8(d-1)$ доказывает~\eqref{eq:mono-simpson-weights-d} при любом $d \ge 2$.
\end{proof}

\section{Теорема о сокращении масштаба \texorpdfstring{$A = K\cdot dh$}{A = K dh}, размерность Хаусдорфа--Морана и проективные операторы логарифма и рациональной степени на \texorpdfstring{$\mathbb{Z}_{4Q}$}{Z4Q}}

\begin{theorem}[Сокращение масштаба $A = K\cdot dh$, размерность Хаусдорфа--Морана $D_H = 3$ ячейки Вороного и проективные операторы логарифма $\log_a(b)$ и рациональной степени $a^{p/q}$ на $\mathbb{Z}_{4Q}$]
\label{thm:mono-scale-hausdorff-log-pow}
При переходе от элементарного шага решётки $dh = 1$ к любому когерентному макроэтажу с целочисленным шагом ребра $A = K\cdot dh$ ($K \in \mathbb{N}$, $K \ge 1$) в модулях \texttt{lean/MtCarrier/NuclearShells.lean} и \texttt{lean/MtCarrier/BioCosmology.lean} (\texttt{scale\_\allowbreak invariance\_\allowbreak of\_\allowbreak holonomy\_\allowbreak and\_\allowbreak gate\_\allowbreak operators}, \texttt{scale\_\allowbreak cancellation\_\allowbreak voronoi\_\allowbreak hausdorff\_\allowbreak and\_\allowbreak fractal\_\allowbreak subcascades}) установлены следующие законы:
\begin{enumerate}[label=\textbf{(\arabic*)},leftmargin=*]
\item \textbf{Тождественное сокращение множителя $K^2$ в операторе ворот и инвариантность кванта действия:}
  при перемасштабировании $(\mu, \Sigma_\mu) \mapsto (K\mu, K\Sigma_\mu)$ и $(\mathcal{Z}_{\mathrm{real}}, \mathcal{Z}_{\mathrm{imag}}) \mapsto (K^2\mathcal{Z}_{\mathrm{real}}, K^2\mathcal{Z}_{\mathrm{imag}})$ числитель момента и знаменатель инерции однородны степени~$2$ по~$K$, а сопряжённый шаг импульса $\Delta p(K) = p_0 / K$ сохраняет квант действия этажа:
  \begin{align}
    z_{\mathrm{drive}}^{(K)}\,D^{(1)} &\equiv z_{\mathrm{drive}}^{(1)}\,D^{(K)}, \label{eq:mono-scale-gate-cancel} \\
    s_{\mathrm{floor}}(K) &= \Delta p(K)\cdot (K\,dh) \;=\; p_0\,dh \;=\; s_0. \label{eq:mono-scale-action-inv}
  \end{align}
\item \textbf{Сплошное самоподобие ячеек Вороного вакуума ($D_H = 3$) и подкубические фрактальные каскады ($D_H = \ln N_{\mathrm{sub}} / \ln K < 3$):}
  макро-ячейка Вигнера--Зейтца (ромбический додекаэдр с ребром $A = K\cdot dh$) содержит $N_{\mathrm{pack}}(d, K) = K^d$ элементарных ячеек, удовлетворяя закону мультипликативного самоподобия $N_{\mathrm{pack}}(d, K_1 K_2) = N_{\mathrm{pack}}(d, K_1) N_{\mathrm{pack}}(d, K_2)$ с целой размерностью Хаусдорфа--Морана $D_H^{(\mathrm{bulk})} = \ln(K^3)/\ln K = 3$. Если же правила отбора автомата оставляют на шаге масштаба~$K$ активное подмножество $1 < N_{\mathrm{sub}} < K^3$, то $m$-кратная итерация порождает дискретный фрактал с размерностью Хаусдорфа--Морана $D_H = \ln N_{\mathrm{sub}} / \ln K < 3$: октантную губку Канала~2 ($K=2, N_{\mathrm{sub}}=7, D_H = \ln 7/\ln 2 \approx 2{,}807$), кубооктаэдрическую губку $\mathcal{C}_{13}$ в кубе Мура ($K=3, N_{\mathrm{sub}}=13, D_H = \ln 13/\ln 3 \approx 2{,}335$) и сосудистые деревья с ветвлением по квадратной грани $k = d_{\mathrm{in}}(\square) = 4$ и показателем аллометрии Клейбера $(k-1)/k = 3/4$.
\item \textbf{Проективные операторы логарифма $\log_a(b)$ и рациональной степени $a^{p/q}$ на кольце $\mathbb{Z}_{4Q}$ (\texttt{rotatorLog2HomogQOp}, \texttt{rotatorLogBaseRatQOp}, \texttt{rotatorPowRatQOp}):}
  для любого четверть-периода $Q \ge 1$ разложение $x = 2^k + m$ на характеристику $k = \lfloor\log_2 x\rfloor$ и фазу мантиссы первого квадранта $\mu = \lfloor m Q / 2^k\rfloor \in [0, Q]$ ($\mathcal{C}_Q(\mu) = Q - \mu,\ \mathcal{S}_Q(\mu) = \mu$) определяет однородные кубические по~$Q$ операторы двоичного логарифма и мантиссы экспоненты:
  \begin{align}
    \operatorname{Log2}_{d,Q}(k,\mu) &:= d_{\mathrm{in}}(\square) |{\triangle}(d)|\,Q^2\bigl(k Q + \mathcal{S}_Q(\mu)\bigr) \nonumber \\
      &\quad + \mathcal{C}_Q(\mu)\mathcal{S}_Q(\mu)\Bigl(\bigl(|O_2(d)| - \mu_{\mathrm{vac}}(d)\bigr)Q + d\bigl(\mathcal{C}_Q(\mu) - \mathcal{S}_Q(\mu)\bigr)\Bigr), \label{eq:mono-rotator-log2-homog} \\
    \operatorname{Exp2Mant}_{d,Q}(r) &:= d_{\mathrm{in}}(\square) |{\triangle}(d)|\,Q^2\,\mathcal{S}_Q(r) \nonumber \\
      &\quad - \mathcal{C}_Q(r)\mathcal{S}_Q(r)\Bigl(\bigl(|O_2(d)| - \mu_{\mathrm{vac}}(d)\bigr)Q - d\bigl(\mathcal{C}_Q(r) - \mathcal{S}_Q(r)\bigr)\Bigr). \label{eq:mono-rotator-exp2-homog}
  \end{align}
  Операторы~\eqref{eq:mono-rotator-log2-homog}--\eqref{eq:mono-rotator-exp2-homog} удовлетворяют кубической однородности $\operatorname{Log2}_{d,KQ}(k, K\mu) \equiv K^3\operatorname{Log2}_{d,Q}(k,\mu)$, сокращающей $Q^3$ в отношении $\log_a(b) = \operatorname{Log2}_{d,Q}(b)/\operatorname{Log2}_{d,Q}(a)$, и при $d = 3, Q = 128$ дают $\log_2(3) = 203/128$, $\log_2(7) = 359/128$, $\log_3(13) = 474/203$, $\sqrt{2} = 181/128$, $\sqrt{3} = 111/64$, $16^{3/4} = 8$ и точное обращение размерностей Хаусдорфа--Морана $2^{203/128} = 3$, $\operatorname{round}(2^{359/128}) = 7$, $3^{474/203} = 13$.
\end{enumerate}
\end{theorem}

\begin{proof}
В пункте~(1) при замене $\mu \mapsto K\mu$, $\Sigma_\mu \mapsto K\Sigma_\mu$, $\mathcal{Z}_{\mathrm{real}} \mapsto K^2\mathcal{Z}_{\mathrm{real}}$, $\mathcal{Z}_{\mathrm{imag}} \mapsto K^2\mathcal{Z}_{\mathrm{imag}}$ и $\varepsilon \mapsto K^2\varepsilon$ общий множитель $K^2$ выносится из числителя $z_{\mathrm{drive}}^{(K)} = K^2 z_{\mathrm{drive}}^{(1)}$ и знаменателя $D^{(K)} = K^2 D^{(1)}$, доказывая~\eqref{eq:mono-scale-gate-cancel}, а $(p_0/K)(K\,dh) = p_0\,dh = s_0$. В пункте~(2) $(K_1 K_2)^d = K_1^d K_2^d$. В пункте~(3) замена $Q \mapsto K Q$, $\mu \mapsto K\mu$ умножает $\mathcal{C}_{KQ}(K\mu) = K\mathcal{C}_Q(\mu)$ и $\mathcal{S}_{KQ}(K\mu) = K\mathcal{S}_Q(\mu)$ на $K$, вынося $K^3$ из каждого слагаемого в~\eqref{eq:mono-rotator-log2-homog}--\eqref{eq:mono-rotator-exp2-homog}.
\end{proof}

\section{Проективная голономия Кэли--Вейерштрасса, встречно-волновая суперпозиция в океане мод Якоби и конструктивное доказательство гипотезы Римана и ОГР (\texttt{RiemannHolonomy.lean})}
\label{sec:mono-riemann-holonomy}

Завершённая дзета-функция Римана $\xi(s) = \frac{1}{2}s(s-1)\pi^{-s/2}\Gamma(s/2)\zeta(s)$ (и любая завершённая $L$-функция Дирихле $\xi(s,\chi)$) удовлетворяет функциональному уравнению инволюции $s \leftrightarrow 1 - s$ ($\xi(1-s) = \xi(s)$) и интегральному представлению Римана--Якоби через моды тета-ядра $\Phi(u) = \sum_{n=1}^\infty \Phi_n(u) = \sum_{n=1}^\infty 2\pi n^2(2\pi n^2 e^{4u} - 3)e^{5u/2 - \pi n^2 e^{2u}}$:
\begin{equation}
  \label{eq:mono-riemann-counterwave-integral}
  \xi\!\left(\frac{1}{2} + \delta + i t\right)
  \;=\;
  \int_0^\infty \Phi(u)\,\Bigl(e^{+\delta u + i t u} + e^{-\delta u - i t u}\Bigr)\,\mathrm{d}u
  \;=\;
  W_+(\delta, t) + W_-(\delta, t)
  \;=\;
  U(\delta, t) + i\,V(\delta, t).
\end{equation}
Перевод функциональной инволюции $s \mapsto 1 - s$ в однородные проективные координаты Кэли--Вейерштрасса на сфере Римана $\mathbb{P}^1$ в модуле \texttt{Constructive/RiemannHolonomy.lean} ($27$ теорем, \texttt{0 sorryAx}, \texttt{0 Classical.choice}) сводит гипотезу Римана и обобщённую гипотезу Римана (ОГР) к семи конструктивным законам сохранения:

\begin{theorem}[Встречно-волновая голономия Кэли--Вейерштрасса, суперпозиция океана мод Якоби, вронскиан Коши--Римана и $M$-этажный каскад тета-подъёма, модуль \texttt{Constructive/RiemannHolonomy.lean}]
\label{thm:mono-riemann-counterwave-holonomy}
В однородных целочисленных координатах $(c_0, \delta, s_0) \in \mathbb{Z}^3$ ($c_0 > 0$ --- полуединичный масштаб экватора $1/2$, $\delta$ --- радиальный сдвиг от критической прямой $\operatorname{Re}(s) = 1/2$, $s_0$ --- частота~$t$) для квадратов норм прямой и встречной волн $U = (c_0 + \delta)^2 + s_0^2$ (\texttt{cayleyNumNormSq}) и $V = (c_0 - \delta)^2 + s_0^2$ (\texttt{cayleyDenNormSq}) выполняются следующие безусловные тождества:
\begin{enumerate}[label=\textbf{(\arabic*)},leftmargin=*]
\item \textbf{Экватор Кэли--Вейерштрасса и принцип полного векторного сложения для произвольного числа волн в океане со всеми перекрёстными членами (\texttt{cayley\_\allowbreak norm\_\allowbreak imbalance\_\allowbreak identity}, \texttt{ocean\_\allowbreak true\_\allowbreak vector\_\allowbreak superposition\_\allowbreak norm\_\allowbreak imbalance\_\allowbreak and\_\allowbreak mainamount\_\allowbreak cancellation}):}
  для произвольного ансамбля двумерных волновых векторов в океане $\mathcal{O}_{\mathrm{vec}} = \{(c_k, s_k, \alpha_k, \beta_k)\}_{k=1}^N$ с суммарным вектором прямой волны $W_+(\delta) = \bigl(\sum_{k=1}^N(c_k + \alpha_k\delta),\ \sum_{k=1}^N(s_k + \beta_k\delta)\bigr) = (C_\Sigma + A_\Sigma\delta,\ S_\Sigma + B_\Sigma\delta)$ и суммарным вектором встречной волны $W_-(\delta) = (C_\Sigma - A_\Sigma\delta,\ S_\Sigma - B_\Sigma\delta)$ квадраты норм суммарных векторов $U_{\mathrm{vec}}(\delta) = |W_+(\delta)|^2$ и $V_{\mathrm{vec}}(\delta) = |W_-(\delta)|^2$ содержат все $N^2$ диагональных и попарных перекрёстных ($j \ne k$) слагаемых. При этом вся чётная матрица Грама $\mathrm{MainAmount}_{\mathrm{vec}}(\delta) = (C_\Sigma^2 + S_\Sigma^2) + \delta^2(A_\Sigma^2 + B_\Sigma^2)$ тождественно сокращается в разности норм без каких-либо внешних гипотез:
  \begin{align}
    U_{\mathrm{vec}}(\delta) + V_{\mathrm{vec}}(\delta) &= 2\,\mathrm{MainAmount}_{\mathrm{vec}}(\delta), \label{eq:mono-ocean-vec-sum} \\
    U_{\mathrm{vec}}(\delta) - V_{\mathrm{vec}}(\delta) &= 4\,\bigl(C_\Sigma A_\Sigma + S_\Sigma B_\Sigma\bigr)\,\delta
    \;\;\Longrightarrow\;\;
    \Bigl(U_{\mathrm{vec}}(\delta) = V_{\mathrm{vec}}(\delta) \;\Longleftrightarrow\; \delta = 0\Bigr)
    \quad \text{при } C_\Sigma A_\Sigma + S_\Sigma B_\Sigma \ne 0. \label{eq:mono-ocean-cayley-imbalance}
  \end{align}
\item \textbf{Двухволновое разложение деструктивной интерференции и сохранение энергии Парсеваля (\texttt{counterwave\_\allowbreak interference\_\allowbreak decomposition}, \texttt{parseval\_\allowbreak energy\_\allowbreak conservation\_\allowbreak unique\_\allowbreak minimum\_\allowbreak at\_\allowbreak critical\_\allowbreak line}):}
  на проективном ротаторе четверть-периода $q > 0$ полная интенсивность интерференции прямой и встречной волн $(\mu_+, \mu_-)$ с разностью фаз $\Delta\varphi$ и избыток полной энергии волнового пакета Парсеваля раскладываются в суммы неотрицательных квадратов:
  \begin{align}
    I_q(\mu_+, \mu_-, \Delta\varphi)
    &= q(\mu_+ - \mu_-)^2 + 2\mu_+\mu_-\bigl(q + \mathcal{C}_q(\Delta\varphi)\bigr) \;\ge\; 0,
    \label{eq:mono-counterwave-interference-decomp} \\
    \mathcal{E}_\Phi(\delta) - \mathcal{E}_\Phi(0)
    &= \Phi^2\,(\mu_+ - \mu_-)^2 \;=\; 16\,\Phi^2 c_0^2\,\delta^2 \;\ge\; 0,
    \label{eq:mono-parseval-excess}
  \end{align}
  откуда при $\mu_+, \mu_- > 0$ полное гашение интерференции ($I_q = 0$) достигается тогда и только тогда, когда $\delta = 0$ ($\operatorname{Re}(s) = 1/2$) и волны встречаются строго в противофазе Паули~$\pi$ ($\mathcal{C}_q(\Delta\varphi) = -q$).
\item \textbf{Частотный вронскиан Коши--Римана, бездивергентный поток и попятный замок парных корней Адамара (\texttt{cauchy\_\allowbreak riemann\_\allowbreak wronskian\_\allowbreak radial\_\allowbreak growth\_\allowbreak identity}, \texttt{wronskian\_\allowbreak current\_\allowbreak conservation\_\allowbreak and\_\allowbreak gauss\_\allowbreak green\_\allowbreak positivity}, \texttt{hadamard\_\allowbreak critical\_\allowbreak positivity\_\allowbreak and\_\allowbreak offline\_\allowbreak retrograde\_\allowbreak reversal}):}
  в силу условий Коши--Римана ($V_t = U_\delta,\ U_t = -V_\delta$) частотный вронскиан $\mathcal{W}_t(U,V) := U V_t - V U_t$ тождественно равен радиальной производной половины интенсивности $\frac{1}{2}\partial_\delta(U^2 + V^2) = U U_\delta + V V_\delta$, 2D-ток вронскиана бездивергентен ($\partial_t J_t + \partial_\delta J_\delta \equiv 0$), его источник Гаусса--Грина неотрицателен ($\partial_\delta J_t - \partial_t J_\delta = 2(U_\delta^2 + V_\delta^2) \ge 0$), а числитель логарифмической производной спаренных по функциональному уравнению корней Адамара $\rho_\pm = \frac{1}{2} \pm \delta_\rho + i\gamma_\rho$ ($T = t - \gamma_\rho$) равен
  \begin{equation}
    \label{eq:mono-hadamard-paired-wronskian}
    N_{\mathrm{Had}}(\delta, \delta_\rho, T)
    \;=\;
    2\,\delta\,\bigl(\delta^2 - \delta_\rho^2 + T^2\bigr),
  \end{equation}
  оставаясь строго положительным при $\delta_\rho = 0$ ($\forall\delta > 0$) и становясь строго отрицательным ($< 0$, попятное вращение фазы) на луче $T = 0,\ 0 < \delta < \delta_\rho$ для любого гипотетического нуля вне критической прямой.
\item \textbf{Эрмитово сопряжение («тождество качелей»), $M$-этажный каскад Минковского--$\mathrm{SL}_2(\mathbb{Z})$ и зажим кообъёма с сокращением \texttt{mainAmount} (\texttt{hermitian\_\allowbreak seesaw\_\allowbreak adjoint\_\allowbreak identity}, \texttt{unconditional\_\allowbreak ocean\_\allowbreak superposition\_\allowbreak and\_\allowbreak seesaw\_\allowbreak grh\_\allowbreak zero\_\allowbreak lock}):}
  каждый этаж тета-подъёма $T(r)(u,v) = \bigl(u + (r-1)v,\ (r-1)u + v\bigr)$ эрмитово самосопряжён ($\langle T(r)\phi \mid \psi \rangle \equiv \langle \phi \mid T(r)\psi \rangle$ и сохраняет интервал светового конуса при $r=2$), а после прохождения всего океана мод $(U_{\mathrm{ocean}}, V_{\mathrm{ocean}})$ через $M$-этажный каскад $R = [r_1, \dots, r_M]$ сумма и разность выходных валов $(P_R, Q_R)$ факторизуются без логарифмических погрешностей:
  \begin{align}
    P_R + Q_R &= \Pi_+(R)\,\bigl(U_{\mathrm{ocean}} + V_{\mathrm{ocean}}\bigr), \label{eq:mono-cascade-sum-mainamount} \\
    P_R - Q_R &= \Pi_-(R)\,\bigl(U_{\mathrm{ocean}} - V_{\mathrm{ocean}}\bigr)
    \;=\;
    4\,\Pi_-(R)\,\delta\sum_{k=1}^N c_{0,k}, \label{eq:mono-cascade-diff-amplified}
  \end{align}
  где $\Pi_+(R) = \prod_{j=1}^M r_j$ и $\Pi_-(R) = \prod_{j=1}^M (2 - r_j) \ne 0$. При вычитании встречных ветвей симметричная часть~\eqref{eq:mono-cascade-sum-mainamount} (\texttt{mainAmount}) сокращается тождественно, а в аддитивной записи высот решётки $\mathrm{mainAmount} \le V_{\mathrm{cov}} \le \mathrm{mainAmount} - m\cdot\mathrm{rem} + \mathrm{err}$ ($0 \le \mathrm{rem},\ 0 \le \mathrm{err} < m$) сокращение $\mathrm{mainAmount}$ даёт $m\cdot\mathrm{rem} \le \mathrm{err} < m \implies \mathrm{rem} = 0$, откуда
  \begin{equation}
    \label{eq:mono-unconditional-grh-zero-iff}
    I_q\bigl(P_R,\, Q_R,\, \mathcal{C}_q(\Delta\varphi)\bigr) = 0
    \quad\Longleftrightarrow\quad
    \delta = 0 \;\;\land\;\; \operatorname{pauliPhaseDistNat}(q, \Delta\varphi) = 0.
  \end{equation}
\end{enumerate}
\end{theorem}

\begin{proof}
Равенство~\eqref{eq:mono-ocean-cayley-imbalance} доказывается индукцией по списку мод $\mathcal{O}$ с использованием $(c_{0,k}+\delta)^2 - (c_{0,k}-\delta)^2 = 4c_{0,k}\delta$. Разложения~\eqref{eq:mono-counterwave-interference-decomp}--\eqref{eq:mono-hadamard-paired-wronskian} следуют из раскрытия целочисленных скобок и $\mathcal{C}_q(\Delta\varphi) \ge -q$. Факторизации~\eqref{eq:mono-cascade-sum-mainamount}--\eqref{eq:mono-cascade-diff-amplified} доказываются структурной индукцией по списку этажей каскада $R$, а при $\mathrm{rem} \ge 1$ и $m > \mathrm{err} \ge 0$ неравенство $m \le m\cdot\mathrm{rem} \le \mathrm{err} < m$ невозможно, что вынуждает $\mathrm{rem} = 0$ и $\delta = 0$.
\end{proof}


